\begin{textsample}{POS Dim 3 – rs – Score 12.00 – rs/rs\_inf\_12.txt}  \label{ex:f3_pos_229}
11.1.
Direitos de Aprendizagem e Desenvolvimento no Campo de Experiências O Eu, o Outro e o Nós CONVIVER com \textbf{crianças} e \textbf{adultos} em \textbf{pequenos} grupos, reconhecendo e respeitando as diferentes identidades e pertencimento étnico-racial, de gênero e religião de seus parceiros. \textbf{BRINCAR} com diferentes parceiros desenvolvendo sua imaginação e solidariedade.
EXPLORAR diferentes formas de interagir com parceiros diversos em situações variadas, ampliando sua noção de mundo e sua sensibilidade em relação aos outros.
PARTICIPAR ativamente das situações do cotidiano, tanto daquelas ligadas ao \textbf{cuidado} de si e do ambiente, como das relativas às atividades propostas pelo/a professor/a.
EXPRESSAR às outras \textbf{crianças} e/ou \textbf{adultos} suas necessidades, emoções, sentimentos, dúvidas, hipóteses, descobertas, opiniões, oposições.
CONHECER -SE e construir uma identidade pessoal e cultural, valorizando suas características e as das outras \textbf{crianças} e \textbf{adultos}, aprendendo a identificar e combater atitudes preconceituosas e discriminatórias. 11.1.1 Objetivos de Aprendizagem e Desenvolvimento Os Objetivos de aprendizagem e desenvolvimento para os \textbf{bebês}, as \textbf{crianças} bem \textbf{pequenas} e as \textbf{crianças} \textbf{pequenas} são assim descritos na BNCC ( 2017 ) e no Referencial Curricular Gaúcho : 11.2 CAMPO DE EXPERIÊNCIAS : CORPO, \textbf{GESTOS} E MOVIMENTOS ( CG ) Na primeira infância, o corpo é o instrumento expressivo e comunicativo que serve de suporte para o desenvolvimento emocional e mental, sendo essencial na construção de \textbf{afetos} e conhecimentos.
Este Campo destaca linguagens que as \textbf{crianças} desde cedo fazem uso e que as orientam em relação ao mundo, como \textbf{gestos}, \textbf{mímicas}, posturas e movimentos expressivos que constituem uma linguagem vital com a qual as \textbf{crianças} expressam emoções, reconhecem sensações, interagem, \textbf{brincam}, ocupam espaços e neles se localizam, construindo conhecimento de si e do mundo.
Ainda neste Campo, destaca -se a importância da construção positiva da autoimagem, como um dos pressupostos básicos para o desenvolvimento integral do sujeito.
A \textbf{criança}, na sua integralidade, conhece e explora o mundo por meio da linguagem corporal, sendo manifestada por meio dos \textbf{gestos}, expressão facial, \textbf{mímicas}, deslocamentos espaciais, manipulação e exploração dos objetos, das \textbf{brincadeiras} e da sua cultura, expressando suas vontades e emoções, vivenciando diferentes experiências em relação ao gênero, à etnia ou raça, à classe, à religião e à sexualidade.
Na Educação \textbf{Infantil}, o corpo é o ponto de partida para as interações e a \textbf{brincadeira}, eixo estruturante da ação pedagógica, e o modo pelo qual as \textbf{crianças} podem se apropriar dos sentidos e funções do mundo social e cultural.
Ao professor cabe assegurar em suas ações pedagógicas o \textbf{cuidado} físico, o desenvolvimento motor, a ampliação de repertório de \textbf{gestos}, o uso do corpo em diferentes espaços promovendo a emancipação e a liberdade, evitando seu cerceamento em situações individuais e coletivas.
O Referencial Curricular Gaúcho, baseado no conceito do Campo de Experiências apresentado na BNCC (2017)⁶, expressa que o movimento, o \textbf{gesto} e o corpo são manifestações de linguagem, através das quais as \textbf{crianças} se expressam e se apropriam de sua corporeidade, na relação com os outros e com o meio.

% matched lemmas: adulto, afeto, bebê, brincadeira, brincar, criança, cuidado, gesto, infantil, mímico, pequeno
\end{textsample}
